\title{LongRoPE: Extending LLM Context Window Beyond 2 Million Tokens}

\begin{abstract}
A substantial context window is a sought-after attribute in large language models (LLMs). Nonetheless, expanding the context length is constrained to approximately 128k tokens because of expensive fine-tuning procedures, limited availability of extensive textual data, and destabilizing effects caused by novel token positions.

This work presents LongRoPE, marking the first instance where the context window of pre-trained LLMs is extended to an impressive \textbf{2048k} tokens. This expansion requires as few as 1k fine-tuning steps, conducted within training sequences of up to 256k tokens, all while preserving the model’s capabilities on its original short context window. The method introduces three primary innovations: (i) we identify and leverage two types of non-uniformities within positional interpolation using an efficient search procedure, resulting in improved initialization for fine-tuning and facilitating an 8$\times$ extension even without additional fine-tuning; (ii) we propose a staged extension approach, in which the model is first fine-tuned at a context length of 256k, after which a second positional interpolation step on the resulting model produces the 2048k-token capacity; (iii) LongRoPE undergoes an additional adjustment at an 8k context length to fully recover performance on short contexts. Comprehensive evaluations with LLaMA2 and Mistral, spanning multiple tasks, validate the benefits of our strategy. LongRoPE can augment models while preserving the original architecture—with only minor updates to the positional embedding layer—and is compatible with nearly all established optimizations. Code will be released at \texttt{https://github.com/microsoft/LongRoPE}

\section{Introduction}

Although Large Language Models (LLMs) have demonstrated impressive capabilities across diverse applications~\cite{gpt4,llama2}, they are frequently constrained by the size of their context window, such as LLaMA2's restriction to 4096 tokens~\cite{llama2}. When input extends beyond this window, the effectiveness of LLMs diminishes due to the model encountering positional contexts for which it has not been trained. This limitation creates difficulties in critical settings, for instance, when performing in-context learning with a large number of exemplars~\cite{Huang2023FewerIM}, or when deploying LLM agents~\cite{park2023generative,madaan2023selfrefine}.

Recent studies indicate that the context window of a pre-trained LLM can be expanded up to approximately 128k tokens through fine-tuning on lengthier documents~\cite{longlora,pi,yarn,zhang2024soaring,liu2023scaling}. However, three primary challenges hinder additional increases in context window size. First, introducing new position indices that lack prior training results in many problematic values, causing out-of-distribution complications that impede convergence during fine-tuning~\cite{pi}. This difficulty is amplified when scaling the context window from 4k to over 1000k, as more than 90\% of positions are unfamiliar to the model. Second, the fine-tuning process typically depends on texts that match the extended window lengths, but current datasets rarely contain such lengthy documents—especially those exceeding 1000k tokens. In addition, processing these extremely long texts demands substantial computational resources, with training times and GPU utilization often becoming impractically high. Third, as the context window becomes exceptionally large, the model’s attention mechanism must allocate its capacity across an enormous number of token positions. This leads to diluted attention and reduced performance for shorter contexts, as previously reported~\cite{pi}.

A commonly employed strategy to address the first challenge involves adapting RoPE positional embeddings~\cite{rope,pi} by scaling novel position indices into the domain of the pretrained embeddings, as depicted in Fig.\ref{fig:overview}. Position Interpolation (PI)~\cite{pi} modifies RoPE's rotary angles using a linear interpolation based on the extension ratio. NTK~\cite{ntk,dynamicntk} proposes non-uniform interpolation and extrapolation across different RoPE dimensions. YaRN~\cite{yarn} further refines this by segmenting RoPE dimensions according to frequency and applying distinct techniques—extrapolation, NTK, or linear interpolation—to each group. Despite these advances, positional embedding inherently contains intricate, non-uniform patterns of information entropy within the Transformer framework. Current interpolation methods fail to fully exploit this complexity, resulting in information degradation and thereby restricting the expansion of context window length.

Section~\ref{sec:analysis} empirically demonstrates two principal observations: \textbf{(1)} Successful positional interpolation requires accounting for two types of non-uniformity—namely, heterogeneous RoPE dimensions and differences across token positions. Interpolation can be applied more conservatively to lower RoPE dimensions and early token locations, although ideal strategies are contingent on the intended extrapolation length. 
\textbf{(2)} Incorporating these non-uniformities during positional interpolation facilitates the preservation of crucial information encoded in the original RoPE, especially with respect to specific dimensions and positional indices. As a result, the degradation due to interpolation is reduced, yielding stronger initialization when adapting models through fine-tuning. Additionally, in settings where fine-tuning is not performed, this approach supports up to an 8$\times$ extension in context length.

Building on these results, we propose \textbf{LongRoPE}, a robust approach designed to increase the context window of LLMs to over 2 \emph{million} tokens. The architecture of LongRoPE incorporates three primary contributions. Firstly, {LongRoPE} leverages multidimensional, non-uniform positional interpolation to its fullest extent. By tailoring rescale factors for the rotation angles within RoPE for each distinct RoPE dimension according to token placement, LongRoPE achieves greater accuracy. Given that the complexity of rescale factor identification grows exponentially relative to the desired extension ratio, {LongRoPE} utilizes an evolutionary search algorithm enhanced by two optimization strategies, significantly improving the efficiency of this process. Fig.~\ref{fig:overview} illustrates an instance of the resulting rescaled RoPE structure.

Subsequently, {LongRoPE} adopts a resource-efficient, staged extension approach to reach a context window of 2048k, thereby avoiding the necessity for fine-tuning directly on extremely lengthy texts, which are both rare and difficult to obtain. The process initiates by identifying a suitable 256k length on the pre-trained LLM and performing fine-tuning at this specific length. Given that our non-uniform positional interpolation supports an 8$\times$ increase in length without additional fine-tuning, we then carry out a second exploration to determine appropriate new RoPE rescale factors using the fine-tuned, length-extended LLM. This methodology culminates in attaining a 2048k context window for both LLaMA2 and Mistral~\cite{mistral}.

In order to address potential declines in performance when reverting to the original, shorter context window, {LongRoPE} maintains its strategy of modifying the RoPE rescale parameters within the expanded LLM model. Utilizing the same methodology employed when increasing from 256k to 2048k context lengths, we also reduce the context window size to 4k and 8k for the LLM that was fine-tuned at 256k, applying our search algorithm to limit positional interpolation. At inference time, when sequences are below 8k tokens, RoPE is adjusted by applying the rescale factors obtained from this search.

A comprehensive series of experiments conducted on multiple LLMs and a range of long-context tasks validate the efficacy of our approach. Our results indicate that LongRoPE consistently sustains low perplexity across evaluation lengths from 4k up to 2048k, attains passkey retrieval accuracy above 90\%, and achieves accuracy on par with standard benchmarks structured within the 4096 context window. Furthermore, LongRoPE is universally compatible with all LLMs utilizing RoPE embeddings. We plan to make both our code and the LongRoPE-2048k models publicly available.

\section{Non-uniformity in Positional Interpolation}
\label{sec:analysis}

\subsection{Preliminary}

Transformer models require explicit positional information, often in the form of position embedding, to represent the order of input tokens. Our work focuses on the RoPE~\cite{rope} position embedding, which is widely used in recent  LLMs. For a token at position index $n$, its corresponding RoPE encoding can be simplified as follows:

\begin{equation}
		\fontsize{7.8}{7.8} \selectfont
	\label{eq:rope}
	\left[ cos(n\theta_0), sin(n\theta_0), cos(n\theta_1), \cdot\cdot\cdot, cos(n\theta_{d/2-1}), sin(n\theta_{d/2-1}) \right]   
\end{equation}

Here, $d$ denotes the embedding dimensionality, 
 $n\theta_i$ corresponds to the rotary angle associated with the token situated at position $n$, and 
$\theta_i=\theta^{-2i/d}$ specifies the set of rotation frequencies. For RoPE, the typical default setting for the base $\theta$ is 10000. 

\textbf{\emph{Context window extension ratio $s$} and positional interpolation}. The parameter $s$ is introduced as the quotient of the extended context length $L'$ by the original context length $L$: $s=\frac{L'}{L}$.

To extend the context window from $L$ to $L'$, current positional interpolation methods suggest downscaling rotation frequencies $\theta_i$ based on extension ratio $s$. Let $\beta=\theta^{2/d}$, and $\lambda$ denote the actual rescale factor related to $s$, we   unify these positional interpolation methods as follows:

\begin{equation}	
	\fontsize{7.5}{7.5} \selectfont
	\label{eq:unified_rope}
	\begin{aligned}
		\left[cos\left(\frac{n}{\lambda(\beta)^0}\right), sin \left(\frac{n}{\lambda(\beta)^0}\right), cos\left(\frac{n}{\lambda(\beta)^1}\right), \cdot\cdot\cdot,  sin \left(\frac{n}{\lambda(\beta)^{d/2-1}}\right) \right]   
	\end{aligned}
\end{equation}

\textbf{Linear positional interpolation (PI)}. The method PI~\cite{pi} involves performing a linear interpolation on position indices, but strictly within the bounds set during pre-training. When extending to a new sequence length with a scaling factor $s$, PI uniformly scales down the rotation angles for every RoPE dimension by a factor of $\lambda=s$. This approach compresses positional information, causing positions to become densely packed and limiting the model’s capability to discern tokens in adjacent positions. As a consequence, PI exhibits reduced efficacy when large extension ratios are applied.

\textbf{NTK-based interpolation and extrapolation}. ~\cite{ntk,dynamicntk} analyze RoPE through the lens of information encoding and employ Neural Tangent Kernel (NTK) theory~\cite{jacot2018neural,tancik2020fourier}. To address the issue of congested positions in PI, they propose redistributing the interpolation load among the various RoPE dimensions. This approach involves applying smaller scaling factors to the lower (high-frequency) dimensions and greater scaling to the higher (low-frequency) ones, thereby enabling both positional interpolation and extrapolation, described by $\lambda=s^{i}$. The enhanced dynamic NTK method~\cite{dynamicntk} further refines this process by modulating the extension ratio for each position according to the current sequence length. While PI typically requires fine-tuning, NTK-informed strategies allow for lengthening the context window even without fine-tuning, though they generally support a maximum extension ratio of 4$\times$.

\textbf{YaRN}~\cite{yarn} achieves notable advancements in positional interpolation by partitioning RoPE dimensions according to frequency into three distinct groups, each processed with a tailored interpolation approach. Dimensions corresponding to high frequencies are treated with extrapolation ($\lambda$=1), low-frequency dimensions utilize linear interpolation (PI), while those in the intermediate frequency range are managed via NTK. The main innovation of YaRN is its frequency-based categorization of RoPE dimensions; however, this grouping relies on empirically-driven, human intervention, which can potentially yield sub-optimal outcomes when applied to novel LLMs.

\subsection{Study on Non-uniform Positional Interpolation}

Building on insights from NTK and YaRN, which derive benefits from introducing non-linearity—particularly through managing distinct frequency behaviors among RoPE dimensions for targeted interpolation and extrapolation—we observe that prevailing approaches to non-linearity are predominantly crafted by hand. This observation leads to two key inquiries: (1) Does existing positional interpolation represent the best possible solution? (2) Might there exist novel, unexamined forms of non-linearity?

In order to address these questions, we employ evolution search (refer to Sec.~\ref{sec:method}) to identify improved non-uniform positional interpolations for LLaMA2-7B. The selection process relies on perplexity scores, utilizing five randomly chosen samples from the PG19~\cite{pg19} validation dataset. Our empirical evaluation leads us to uncover several important insights.

\textbf{Finding 1}: \textit{RoPE dimensions exhibit substantial non-uniformities, which are not effectively handled by current positional interpolation methods.}

We determine the optimal value of $\lambda$ for each RoPE dimension as defined in Eq.~\ref{eq:unified_rope}. Table~\ref{tbl:exp1} presents a perplexity comparison for LLaMA2-7B across various methods on the PG19 and Proof-pile~\cite{proof-pile} test datasets, evaluated without fine-tuning. The results from our search exhibit marked gains, indicating that commonly used linear (PI) and uneven (Dynamic-NTK and YaRN) interpolations do not yield the best performance. In particular, YaRN lags behind both PI and NTK on the PG19 dataset, as it fails to achieve the intended context window length for LLMs that have not been fine-tuned. For instance, with an 8k context size, YaRN’s perplexity rises sharply past the 7k point.

In our investigation, the rescaling factors $\lambda$ in Eq.~\ref{eq:unified_rope} are found to be non-uniform, contrasting with the constant scaling factor $s$ utilized in the approaches of PI, NTK's formula, and YaRN's group-wise method. Employing these non-uniform factors leads to notable enhancements in the language modeling performance of LLaMA2—measured via perplexity—for extended context windows of 8k and 16k, without the need for fine-tuning. This improvement stems from the fact that the resulting positional embedding more faithfully maintains the structure of the original RoPE, particularly in critical dimensions, thereby lessening the model's challenge in differentiating between closely positioned tokens.

\textbf{Finding 2}: \textit{RoPE for the initial tokens in the input sequence should be extrapolated with less interpolation.} 

We posit that the RoPE applied to the first $\hat{n}$ tokens of input sequences should involve minimal interpolation, since these tokens often obtain significant attention weights and play a pivotal role in the attention mechanism, as supported by findings in Streaming LLM~\cite{streamingllm} and LM-Infinite~\cite{han2023lminfinite}. To investigate this, we experiment with context window extensions up to 8k and 16k using PI and NTK, while withholding interpolation for the initial $\hat n$ tokens (where $\hat n$ ranges over 0, 2, ..., 256). If $\hat n$=0, the procedure defaults to standard PI and NTK. Table~\ref{tbl:tokenposition} yields two primary conclusions: \textbf{(1)} excluding position interpolation for the earliest tokens leads to demonstrable gains in performance; \textbf{(2)} the ideal choice of $\hat n$ is contingent upon the specific length of the context window being extended.

\textbf{Finding 3}: \textit{Non-uniform positional interpolation effectively extends LLM context window in both fine-tuning and non-fine-tuning settings.} 

Although our optimized non-uniform position interpolation enhances extension performance to 8k and 16k contexts without additional training, extending further necessitates fine-tuning. Therefore, we adapt LLaMA2-7B using our discovered RoPE scheme for context windows up to 64k tokens (refer to Appendix for detailed configurations). As demonstrated in Table~\ref{tbl:finetune}, our approach markedly surpasses both PI and YaRN, regardless of whether fine-tuning has been applied to LLaMA2-7B. This superiority arises from the judicious application of non-uniform positional interpolation, which reduces information degradation and enables more effective initialization when fine-tuning.

\textbf{Summary}. This work identifies two types of non-uniformity: differences among RoPE dimensions and across token positions. Strategic exploitation of these non-uniformities in positional interpolation yields substantial gains in large language model context extension performance.

\section{LongRoPE}

\label{sec:method}
Building upon these observations, we propose LongRoPE, which incorporates an optimized search method to harness both types of non-uniformity. This approach enables the extension of the LLM context window to surpass 2 million tokens.

\subsection{Problem Formulation}

These two forms of non-uniformity can greatly expand the set of possible solutions and complicate the process of optimization. To manage this challenge, we reformulate the multidimensional non-uniform position interpolation optimization task into a search problem.

For a LLM targeting a  context window size of $L'$ and lengthy input documents $\mathbf{X}$, where each $\mathbf{x}\in\mathbf{X}$ surpasses $L'$ in token length, we denote the original rotary angle of the $i^{th}$ dimension in RoPE embedding at token position $n$ as  $\frac{n}{\beta^i}$. The optimization problem is then formulated as follows:

\begin{equation}
\begin{gathered}
		\label{eq:problem}

\underset {\mathbf{x}\in\mathbf{X}; \, |\mathbf{x}|\ge L'}{ \operatorname {arg\,min} } \,  \mathcal{L}\, \left( \text{LLM} (\text{RoPE}, \mathbf{X})\right),	\text{with}
\\
\scalebox{0.93}{
$\underset {\begin{subarray}{l} i=0,\cdots,\frac{d}{2}-1;\\ n\in[ 0,|\mathbf{x}|);\end{subarray}}{\text{RoPE}_(n)}= {\left[\cdots,\cos\left(\mathbb{I}(\hat{\lambda}_{i}, \hat{n})\frac{n}{\beta^i}\right),  \sin\left(\mathbb{I}(\hat{\lambda}_{i}, \hat{n})\frac{n}{\beta^i}\right),\cdots\right] }$}\\
	\text{and} \; \mathbb{I}(\hat{\lambda}_{i},\hat{n})=\begin{cases}
	1&\mbox{\;  $n<\hat{n}$}\\
{\frac{1}{\lambda}_{i}}&\mbox{\; $n\geq\hat{n}$}
\end{cases}
	\end{gathered}
\end{equation}

We define a collection of rescale factors, $\mathbb{I}(\hat{\lambda}_{i},\hat{n})$, which serves to address both identified types of non-uniformities. Here, $\hat{\lambda_i}$ corresponds to non-uniformity across RoPE dimensions, while $\hat{n}$ pertains to irregularities related to token positions. In practice, $\mathbb{I}(\hat{\lambda}_{i},\hat{n})$ is employed to adjust the rotation angle within the $i^{th}$ RoPE dimension, with $\hat\lambda_{i}$ acting as the rescale factor and $\hat{n}$ representing a threshold for token position. For token positions less than $\hat{n}$, the adjustment via $\hat\lambda_{i}$ is not applied, so the standard RoPE rotary angle $\frac{n}{\beta^i}$ is preserved. When token positions satisfy $n\ge\hat{n}$, the rescale factor becomes effective.

Suppose we specify a desired context window length $L'$. The task then is to determine the set of rescaling factors—namely, $\mathbb{I}(\hat{\lambda}_{0},\hat{n})$, $\mathbb{I}(\hat{\lambda}_{1},\hat{n})$, $\ldots$, $\mathbb{I}(\hat{\lambda}_{i},\hat{n})$, through the $d^{th}$ RoPE component—which are optimal for this context. By rescaling the $\text{RoPE}$ in the target $\text{LLM}$ accordingly, we seek to minimize the loss for predicting the next token, denoted by $\mathcal{L}$ (perplexity), when processing input sequences $\mathbf{X}$ with a length of $L'$.

\subsection{Searching the Non-uniform Position Interpolation}

To address the challenge presented in Eq.~\ref{eq:problem}, we propose a straightforward but powerful approach that identifies the optimal RoPE rescaling parameters, thereby leveraging the diverse non-uniform characteristics inherent in position embeddings across multiple dimensions.

\textbf{Search space}. We construct an expansive search space to encompass both non-uniform characteristics under consideration. Table~\ref{tbl:searchspace} provides an overview of this search space. Rather than searching the indicator function $\mathbb{I}(\hat{\lambda}_{i},\hat{n})$ directly, we opt for a simplified approach where the search targets $\lambda_i$ and $\hat{n}$, with $\hat{\lambda}_{i}$ being defined as $1/\lambda_i$. The search space supports determining a unique rescaling factor for each dimension of the RoPE embedding. According to Table~\ref{tbl:searchspace}, $\lambda_i$ can be selected from values ranging between 1.0 (representing straightforward extrapolation) and $s\times1.25$ (allowing for broader interpolation than PI), incrementing in steps of 0.01. Here, $s$ represents the ratio by which the target context window is extended.

The parameter $\hat{n}$ determines how many initial token positions preserve their original RoPE embedding, without applying position interpolation. In practice, we let $\hat{n}$ vary among the set \{0, 1, 2, 4, 8, 12, 16, 20, 24, 28, 32, 64, 128, 256\} during the search. Setting $\hat{n}=0$ results in all token positions relying solely on the optimized rescale factors.

\textbf{Evolution-based search}. As shown in Table~\ref{tbl:searchspace}, the breadth of our search space includes a vast array of positional interpolation configurations, which creates considerable difficulty for thorough and efficient search. For instance, when employing an extension of $s=4\times$, one faces $400^{128/2}\times14$ = $4\times10^{167}$ possible candidates. This exponential growth of the search space becomes even more pronounced at higher extension ratios. To mitigate this issue, we implement evolution search~\cite{guo2020single} complemented by two optimization strategies that substantially enhance search speed. The overarching search process is detailed in Algorithm~\ref{alg:search}.

\textit{Optimized initial population generation}. Rather than selecting all $P$ rescale factors at random for the initial population, we explicitly include the three RoPE rescale factors associated with PI, NTK, and YaRN among the starting individuals. The rest of the $P$-3 individuals are produced by applying random mutations to these three rescale factors with a probability $p$.

\textit{Monotonically non-decreasing constraint}. Once the initial population has been produced, we evaluate the LLM perplexity for every candidate. This involves assigning the relevant RoPE rescale factors to the target LLM and measuring the perplexity over the input $\mathbf{X}$. The individuals with the best perplexity scores (top-k) are selected as parents for subsequent evolutionary steps. Nonetheless, due to the expansive nature of the search space, unsophisticated mutation and crossover procedures can lead to the selection of suboptimal candidates, resulting in a surplus of perplexity evaluations. This inefficiency is exacerbated when $L'$ is large, as each perplexity computation necessitates expensive inference.

To tackle this issue, we enforce a monotonicity restriction such that the rescaled RoPE factors sampled fulfill $\lambda_i \le \lambda_{i+1}$ for all $i$. Only those RoPE configurations that adhere to this non-decreasing property are utilized during LLM perplexity assessment, thereby greatly curtailing the computational expense of the search process. More concretely, we stipulate that $\lambda_i$ must grow in a monotonic fashion with respect to the RoPE dimension index ($i$ ranging from 0 to 63). This requirement is inspired by the NTK theory~\cite{jacot2018neural,tancik2020fourier,ntk}, which proposes that for dimensions with higher frequencies (i.e., smaller $i$), less interpolation is needed (implying smaller $\lambda_i$), while dimensions with lower frequencies (i.e., larger $i$) permit greater interpolation (hence, larger $\lambda_i$).

\begin{algorithm}[t]
	\small
	\caption{The search algorithm}
	\textbf{Input:} target LLM, input samples $\mathbf{X}$, population size $P$, mutation size $N_1$, crossover size $N_2$, maximum iterations $\mathcal{T}$, mutation probability $p$\\

	\begin{algorithmic}[1]
		\label{alg:search}
		\STATE Top-k $\gets \phi$;	
		\STATE $\text{P}_0$ $\gets$ \textit{Initialize\_population\_with\_optimization}($P$, $\mathbf{X}$, $p$);
		\FOR{$i = 1$ \textbf{to} $\mathcal{T}$}
		\STATE \textit{Compute\_perplexity}(LLM, $\text{P}_{i-1}$, $\mathbf{X}$);
		\STATE Top-k $\gets$ \textit{Update\_Topk}(Top-k, $\text{P}_{i-1}$);
		\STATE $\text{P}_{mutation}$ $\gets$ \textit{Mutation\_with\_mono\_constraint}(Top-k, $N_1$, $p$);
		\STATE $\text{P}_{crossover}$ $\gets$ \textit{Crossover\_with\_mono\_constraint}(Top-k, $N_2$);
		\STATE $\text{P}_i$ $\gets$ $\text{P}_{mutation} \cup \text{P}_{crossover} \cup$ Top-k;
		\ENDFOR
		\STATE Output the individual from Top-k exhibiting the minimum perplexity;
	\end{algorithmic}
\end{algorithm}

\textbf{8$\times$ extension without fine-tuning}. Leveraging our evolutionary search approach, we successfully determine non-uniform RoPE rescale factors that selectively retain critical dimensions and positions, thereby reducing the loss of information due to interpolation. As shown in Fig.\ref{fig:non-ft}, this strategy enables us to enlarge LLaMA2's context window from 4k up to 32k tokens without the need for fine-tuning. By comparison, alternative techniques like PI and both non-uniform NTK and YaRN result in a significant increase in perplexity beyond a 2$\times$ context length extension.

\subsection{Extending LLM Context Window to 2048K}
\label{sec:extendto2048k}

\textbf{Progressive extension to 2048k}. In this section, we present our approach for increasing the context window of pre-trained LLMs from the conventional 4k up to 2048k. Our non-uniform positional interpolation technique is capable of achieving up to an 8$\times$ context expansion without the need for fine-tuning. However, when attempting to extend to a much larger window size (such as 512$\times$), fine-tuning becomes indispensable. One strategy involves determining RoPE rescaling factors appropriate for the 2048k window and performing fine-tuning accordingly. Despite this, such an approach is impeded by the significant computational and financial costs of training at these scales. Additionally, our observations suggest that achieving effective fine-tuning for LLMs at such high extension ratios is quite difficult (refer to Appendix for details).

Luckily, LongRoPE demonstrates strong performance with both the base model and extended versions of LLMs following fine-tuning. To this end, we propose an efficient and incremental approach that reaches the desired 2048k context length using only 1k fine-tuning iterations, with training sequences capped at 256k tokens.

$\Diamond$ \textit{LongRoPE search for scaling pre-trained LLMs to 256k}. Using LLaMA2 as a case study, we systematically explore context window expansion to 128k and 256k tokens. At this phase, the scaling factors reach 32$\times$ and 64$\times$ for 128k and 256k contexts, respectively.

$\Diamond$ \textit{Fine-tuning to 256k}. Following pre-training, we adjust the LLM to support a 256k context window via fine-tuning. In detail, we initiate fine-tuning of LLaMA2 for 400 steps with RoPE rescaling set for 128k. Once this stage is complete, we switch the RoPE rescale factors to 256k on the obtained checkpoint, continuing fine-tuning for an additional 600 steps. This approach demonstrates greater efficiency compared to attempting a direct fine-tuning process to 256k.

$\Diamond$ \textit{Utilizing LongRoPE to expand fine-tuned extended LLM context to 2048k}. In the last phase, we conduct an additional search on the LLM previously fine-tuned for sequences of 256k length. Through this process, the model achieves an exceptionally vast context window of 2048k, requiring no extra fine-tuning. The extension ratio ultimately reaches $512\times$.

\textbf{Restoring performance in the original context window}.  Upon scaling up the context window to a massive 2048k tokens, we observe diminished model effectiveness within the initial, shorter window. This reduction in performance has been documented as a typical consequence of positional interpolation~\cite{pi}, which restricts the range available for position embeddings at lower token positions, confining them to a relatively tight segment and consequently hampering the language model’s capabilities. When employing an extension factor of 512$\times$, the positional representations for tokens in the original 4k window are densely packed, exacerbating this issue.

To address this issue, we conduct an additional evolution search on the expanded LLM, specifically tuning the RoPE rescale parameters for shorter context lengths, such as 4k and 8k tokens. For these shorter contexts, we lower the upper bound of the search space for $\lambda$, as the need for positional interpolation diminishes. At inference time, the LLM adaptively modifies the relevant RoPE rescale values.

\section{LongRoPE}

\label{sec:method}
Building upon these insights, we propose LongRoPE, which initially employs an optimized search procedure to take full advantage of the identified non-uniform characteristics, and subsequently leverages this mechanism to expand the context window of LLMs to surpass 2 million tokens.

\subsection{Problem Formulation}

These two forms of non-uniformity substantially expand the set of potential solutions and increase the complexity of the optimization process. To tackle this challenge, we recast the multidimensional non-uniform position interpolation optimization as a search-based problem.

For a LLM targeting a  context window size of $L'$ and lengthy input documents $\mathbf{X}$, where each $\mathbf{x}\in\mathbf{X}$ surpasses $L'$ in token length, we denote the original rotary angle of the $i^{th}$ dimension in RoPE embedding at token position $n$ as  $\frac{n}{\beta^i}$. The optimization problem is then formulated as follows:

\begin{equation}
\begin{gathered}
		\label{eq:problem}

\underset {\mathbf{x}\in\mathbf{X}; \, |\mathbf{x}|\ge L'}{ \operatorname {arg\,min} } \,  \mathcal{L}\, \left( \text{LLM} (\text{RoPE}, \mathbf{X})\right),	\text{with \;} 
\\
\scalebox{0.93}{
$\underset {\begin{subarray}{l} i=0,\cdot\cdot,\frac{d}{2}-1;\\ n\in[ 0,|\mathbf{x}|);\end{subarray}}{\text{RoPE}_(n)}= {\left[\ldots,cos\left(\mathbb{I}(\hat{\lambda}_{i}, \hat{n})\times\frac{n}{\beta^i}\right),  sin\left(\mathbb{I}(\hat{\lambda}_{i}, \hat{n})\times\frac{n}{\beta^i}\right),\ldots\right] }$}\\
	\text{such that \;} \mathbb{I}(\hat{\lambda}_{i},\hat{n})=\begin{cases}
	1&\text{\;  if $n<\hat{n}$}\\
{\frac{1}{\lambda}_{i}}&\text{\; if $n\geq\hat{n}$}
\end{cases}
	\end{gathered}
\end{equation}

We define a set of rescale factors, $\mathbb{I}(\hat{\lambda}_{i},\hat{n})$, to address both types of non-uniformities. Here, $\hat{\lambda_i}$ captures non-uniformity across RoPE dimensions, and $\hat{n}$ pertains to irregularities in token positions. The function $\mathbb{I}(\hat{\lambda}_{i},\hat{n})$ is specifically employed to adjust the rotation angle for the $i^{th}$ RoPE dimension, with $\hat\lambda_{i}$ acting as the adjustment coefficient and $\hat{n}$ serving as the position cutoff. For token positions earlier than $\hat{n}$ (i.e., the first $\hat{n}$-1 positions), the adjustment $\hat\lambda_{i}$ is omitted, preserving the original RoPE rotation angle, given by $\frac{n}{\beta^i}$. When $n\ge\hat{n}$, however, the rescale factor is introduced and applied accordingly.

Assuming a desired context window size of $L'$, the task is to determine the best rescale factors—$\mathbb{I}(\hat{\lambda}_{0},\hat{n})$, $\mathbb{I}(\hat{\lambda}_{1},\hat{n})$, ..., $\mathbb{I}(\hat{\lambda}_{i},\hat{n})$—for each RoPE dimension from $1^{st}$ through $d^{th}$. Through this approach, the modified $\text{LLM}$, utilizing the rescaled $\text{RoPE}$, is expected to reach the lowest possible next token prediction loss $\mathcal{L}$ (that is, perplexity) when processing input samples $\mathbf{X}$ whose token lengths are $L'$.

\subsection{Searching the Non-uniform Position Interpolation}

To address the challenge outlined in Eq.~\ref{eq:problem}, we present a straightforward but powerful approach that aims to identify the ideal RoPE rescale factors, enabling comprehensive utilization of the multidimensional variations inherent in position embeddings.

\textbf{Search space}. We construct an extensive search space to account for both forms of non-uniformity. The details of this search space are presented in Table~\ref{tbl:searchspace}. In particular, our approach permits the optimization of an individual rescale parameter for every dimension within the RoPE embedding. For clarity and ease of search space formulation, we optimize over $\lambda_i$ and $\hat{n}$, as opposed to directly searching for $\mathbb{I}(\hat{\lambda}_{i},\hat{n})$, where $\hat{\lambda}_{i}$ is defined as $1/\lambda_i$. According to Table~\ref{tbl:searchspace}, the allowed range for $\lambda_i$ starts at 1.0 (which corresponds to pure extrapolation) and extends up to $s\times1.25$ (representing more aggressive interpolation than PI), with increments of 0.01. Here, $s$ denotes the intended ratio for extending the context window.

$\hat{n}$ determines how many of the first token positions are preserved with their original positional encoding, specifically maintaining RoPE embeddings without interpolation. Based on experimental settings, $\hat{n}$ is explored over the set \{0, 1, 2, 4, 8, 12, 16, 20, 24, 28, 32, 64, 128, 256\}. In cases where $\hat{n}=0$, position interpolation is applied to every token position, utilizing the optimized rescale factors throughout.

\textbf{Evolution-based search}. The search space described in Table~\ref{tbl:searchspace} includes a vast array of positional interpolation approaches, making comprehensive exploration highly demanding. For instance, increasing the extension to $s=4\times$ results in $400^{128/2}\times14$=4$\times10^{167}$ possible combinations. As the extension ratio grows, the search space proliferates at an exponential rate. To manage this complexity, we adopt evolution search~\cite{guo2020single} and implement two optimization strategies designed to significantly enhance search efficiency. The full search process is presented in Algorithm~\ref{alg:search}.

\textit{Optimized initial population generation}. Rather than relying solely on random initialization for the population of $P$ rescale factors, we incorporate the three RoPE rescale factors associated with PI, NTK, and YaRN directly as individuals in the starting population. The other $P-3$ members of the population are created by randomly mutating the aforementioned three rescale factors, applying mutations with probability $p$.

\textit{Monotonically non-decreasing constraint}. Once the initial population has been produced, we evaluate LLM perplexity for every candidate. This involves assigning the designated RoPE rescale factors to the relevant LLM and then determining the perplexity value for the input $\mathbf{X}$. The best-performing $k$ candidates are chosen as parents for the subsequent evolutionary steps. Nevertheless, due to the immense size of the search space, unstructured application of mutation and crossover operators may frequently result in suboptimal candidates, thus requiring many extraneous perplexity evaluations. This inefficiency is particularly pronounced when $L'$ assumes a large value, because each perplexity calculation involves computationally expensive inference.

To mitigate this issue, a constraint enforcing non-decreasing monotonicity is placed on the sequence of rescaled RoPE factors: $\lambda_i\le \lambda_{i+1}$. Only RoPE variants meeting this requirement are utilized within the LLM during perplexity assessments, which sharply curtails computational expense by narrowing the search space. Explicitly, we mandate that $\lambda_i$ strictly grows with respect to the RoPE dimension index ($i = 0, \ldots, 63$). This monotonic relationship across dimensions aligns with NTK theory~\cite{jacot2018neural,tancik2020fourier,ntk}, which argues that dimensions associated with higher frequencies benefit from lower interpolation—hence smaller values for $\lambda_i$—whereas dimensions corresponding to lower frequencies accommodate greater interpolation, thus warranting larger $\lambda_i$ values.

\begin{algorithm}[t]
	\small
	\caption{The search algorithm}
	\textbf{Input:} target LLM, input samples $\mathbf{X}$, population size $P$, mutation size $N_1$, crossover size $N_2$, max iterations $\mathcal{T}$, mutate probability $p$\\

	\begin{algorithmic}[1]
		\label{alg:search}
		\STATE Initialize Top-k to $\phi$;	
		\STATE Set initial population $\text{P}_0$ using \textit{Initialize\_population\_with\_optimization} ($P$, $\mathbf{X}$, $p$); 
		\FOR{$i$ from $1$ to $\mathcal{T}$}
		\STATE Execute \textit{Compute\_perplexity} (LLM, $\text{P}_{i-1}$, $\mathbf{X}$);
		\STATE Update Top-k using \textit{Update\_Topk} (Top-k, $\text{P}_{i-1}$);
		\STATE Generate mutation set $\text{P}_{mutation}$ via \textit{Mutation\_with\_mono\_constraint} (Top-k, $N_1$, $p$); 
		\STATE Obtain crossover set $\text{P}_{crossover}$ through \textit{Crossover\_with\_mono\_constraint} (Top-k, $N_2$);
		\STATE Form new population $\text{P}_i$ as union of $\text{P}_{mutation}$, $\text{P}_{crossover}$, and Top-k;
		\ENDFOR
		\STATE Output the individual with the minimum perplexity from Top-k;
	\end{algorithmic}
\end{algorithm}

\textbf{8$\times$ extension achieved without fine-tuning}. The evolutionary search algorithm successfully determines non-uniform RoPE rescale coefficients that retain critical dimension and position features, thereby reducing information degradation associated with interpolation. As shown in Fig.\ref{fig:non-ft}, this technique enables LLaMA2’s context window to expand from 4k up to 32k with no need for fine-tuning. Conversely, alternative approaches like PI, and non-uniform NTK and YaRN result in significant increases in perplexity when attempting extensions beyond 2$\times$.

\subsection{Extending LLM Context Window to 2048K}
\label{sec:extendto2048k}

\textbf{Progressive extension to 2048k}. Here, we present our approach for scaling the context window of pre-trained LLMs from the standard 4k to beyond 2048k. Our non-uniform positional interpolation method has already demonstrated the ability to realize an 8$\times$ expansion without necessitating fine-tuning. When seeking even greater expansions, such as up to 512$\times$, fine-tuning becomes essential. A common technique involves identifying the RoPE rescaling factors that correspond to the target 2048k context size and subsequently applying fine-tuning. Nevertheless, this approach is limited by the extremely high computational costs associated with training. Additionally, our observations indicate that performing high-quality fine-tuning under such large extension ratios is difficult (refer to Appendix).

Luckily, LongRoPE demonstrates efficacy when applied to both the baseline and the length-augmented LLM. To this end, we propose a streamlined and incremental strategy, enabling the model to reach the intended 2048k context length with only 1k steps of fine-tuning, while keeping the training sequence length capped at 256k.

$\Diamond$ \textit{LongRoPE search for scaling pre-trained LLMs to a 256k context}. Using LLaMA2 as a case study, we perform a search process aiming for enlarged context window sizes of 128k and 256k. The corresponding extension factors during this phase are 32$\times$ and 64$\times$, respectively.

$\Diamond$ \textit{Fine-tuning to 256k}. In this stage, the pre-trained LLM is adapted to handle a context window of 256k tokens. The process involves an initial fine-tuning of LLaMA2 over 400 steps utilizing RoPE rescaling parameters suited for 128k tokens. Upon completion, the RoPE rescaled parameters are updated to accommodate 256k tokens, and an extra 600 steps of fine-tuning are carried out. This sequential approach demonstrates greater efficiency compared to immediately fine-tuning for 256k tokens.

$\Diamond$ \textit{Increasing fine-tuned extended LLM capacity to 2048k using LongRoPE search}. In the last step, we carry out an additional search using the previously fine-tuned LLM with a 256k context length. This leads to an exceptionally expanded context window of 2048k, achieved without any subsequent fine-tuning. The overall expansion factor is $512\times$.

\textbf{Recovering performance in shorter context windows}. When expanding the context window to a massive 2048k, we observe a degradation in performance inside the initial context window span. This phenomenon has been documented in relation to positional interpolation~\cite{pi}, which compresses the position embeddings in higher dimensions into a limited space within the original context window, thereby impairing the language model’s efficacy. Specifically, using an extension ratio of 512$\times$ results in considerable positional crowding for tokens within the first 4k context window.

To address this issue, we conduct an additional evolution search on the expanded LLM, fine-tuning the RoPE rescale factors specifically for short context windows (such as 4k and 8k). Because shorter contexts demand minimal positional interpolation, we decrease the upper limit for $\lambda$ considered during search. In inference, the LLM adaptively modifies the relevant RoPE rescale parameters as needed.

 \section{Experiments}

\subsection{Setup}
\textbf{Evaluation Tasks and Models}. LongRoPE is integrated with LLaMA2-7B and Mistral-7B and assessed across three dimensions: (1) perplexity scores for LLMs handling extended-contexts on lengthy documents; (2) performance on the Passkey retrieval task, which tests the model's capability to locate a passkey amid extensive distractor content; and (3) outcomes on conventional LLM benchmarks with a confined context length of 4096 tokens.

\textbf{Fine-tuning}. For LLaMA2, our procedure employs a learning rate of $2\times 10^{-5}$ with linear decay across a global batch size of 32. The initial phase consists of 400 fine-tuning steps using the Redpajama~\cite{redpajama} dataset, divided into segments of 128k tokens each and delimited by BOS and EOS tokens. Building upon the resulting checkpoint, a further 600 training steps are performed to extend the context window to 256k. The 128k context configuration utilizes 8 A100 GPUs within the distributed framework~\cite{cube}, while the 256k variant demands 16 A100 GPUs. For the Mistral model, we set a fixed learning rate at $1\times 10^{-6}$ with a global batch size of 64. Both the 128k and 256k context versions adhere to the protocol described in YaRN~\cite{yarn}, performing 400 steps on Together Computer’s Long-Data Collections~\cite{mistraldata} with a sequence length of 16k. Training is conducted using 4 A100 GPUs.

\textbf{Search}. When considering a target window size up to 256k, parameters are set as follows: $P$=64, $N_1$=$N_2$=16, $p$=0.3, $\mathcal{T}$=40. In each iteration, the top-32 individuals are chosen for mutation and crossover. Perplexity evaluation is performed on 5 randomly selected validation samples from PG19, each meeting or exceeding the desired context length. For window sizes greater than 512k, we reduce the population, as well as mutation and crossover counts, by 50%. In this regime, perplexity is evaluated using 3 randomly sampled passages from the Pile-Books3~\cite{pile} validation set.

\textbf{Baselines}. To obtain a 2048k context length, we conducted fine-tuning procedures utilizing models pre-trained with 128k and 256k window sizes. Consequently, this produces LongRoPE-2048k (ft=128k) for LLaMA2 and LongRoPE-2048k (ft=256k) for Mistral. The four variants are evaluated against leading context window extension baselines, focusing on open-source large language models that have undergone post-interpolation fine-tuning via PI, NTK, and YaRN methods. Representative models in this category include Together-32k~\cite{together}, Code LLaMA~\cite{codellama}, LongLoRA-full-FT-100k~\cite{longlora}, as well as YaRN-LLaMA and YaRN-Mistral~\cite{yarn}.

\subsection{Main Results}

\label{sec:mainresult}
\textbf{Language modeling with long sequences up to 256k tokens}. Our analysis starts with a performance comparison to advanced long-context language models at an evaluation length of 256k tokens. To illustrate the broad applicability of our approach, we employ two benchmark datasets: the test splits of Proof-pile~\cite{pg19} and PG19~\cite{pile}. Perplexity measurements are recorded at multiple context lengths using a sliding window of size 256. For PG19, the evaluation covers the complete test split, comprising 100 documents. For Proof-pile, following the sampling protocol of YaRN~\cite{yarn}, we select 10 random samples, ensuring each has a minimum length of 128k.

Table~\ref{tbl:proof-pile} and Table~\ref{tbl:pg19} present a perplexity comparison between LLaMA2 and Mistral, both enhanced with various interpolation strategies, on the Proof-pile and PG19 datasets. Two main findings are notable: \textbf{(1)} the models augmented by our methods display a consistent decrease in perplexity as the evaluation length increases from 4k up to 256k, demonstrating their capacity to utilize extended context windows effectively. \textbf{(2)} Remarkably, even when handling a context window that is 16$\times$ larger—a scenario that usually hampers short-context performance—our LongRoPE-2048k variants achieve superior perplexity to leading baselines within the 256k context range.

\textbf{Modeling extended language sequences surpassing 2000k}. To assess performance on very lengthy texts, we utilize the Books3~\cite{pile} dataset. For efficient evaluation, 20 books longer than 2048k are randomly chosen, and a sliding window of 256k is applied.

Referencing Table~\ref{tbl:books3}, LongRoPE demonstrates effectiveness in scaling up the context window of both LLaMA2-7B and Mistral-7B to 2048k tokens. Within more moderate context lengths ranging from 8k to 128k, LongRoPE exhibits perplexity levels that are equal to or surpass those of the existing baselines. Distinct performance trends are evident between LLaMA2 and Mistral under the 2048k configuration. Mistral shows advantageous results at shorter context lengths; however, its perplexity rises above 7 past the 256k mark. LLaMA2, consistent with prior expectations, maintains a declining perplexity as the context window expands, though slight increases are noticed at the 1024k and 2048k points. Notably, LongRoPE-2048k achieves improved outcomes on LLaMA2 when fine-tuned at 256k as opposed to 128k, attributable to a reduced secondary extension ratio (specifically, 8$\times$ compared to 16$\times$). Conversely, Mistral attains optimal performance at a 128k fine-tuning window. This discrepancy stems primarily from the adoption of YaRN's protocol, in which both 128k and 256k fine-tuning on Mistral employ a training length of 16k, thus constraining Mistral’s subsequent capability to further increase the context window following fine-tuning.

\textbf{Passkey retrieval}. Next, we examine the impact of context window size on generation performance by employing a synthetic passkey retrieval task, as introduced by~\cite{passkey}. This setup requires the model to identify and extract a randomly inserted five-digit number (the passkey) embedded within a lengthy document. The precise structure of the prompt is described in the appendix. For this evaluation, we conduct 10 runs of the task, each time inserting the passkey at a randomly selected position that is uniformly distributed over the evaluation context window.

Figure~\ref{fig:passkey} presents a comparison of retrieval accuracy against baseline methods. The accuracy of existing approaches declines sharply to zero once the sequence length exceeds 128k. In sharp contrast, LongRoPE-LLaMA2-2048k (ft=256k) demonstrates robust retrieval performance—exceeding 90\% accuracy—throughout the entire range from 4k to 2048k, even under the challenging conditions of extracting a passkey from a dataset with millions of tokens. Meanwhile, LongRoPE-Mistral-2048k (ft=128k) maintains perfect accuracy up to 1800k tokens, with performance dropping to 60\% at 2048k. This drop is consistent with the findings in Table~\ref{tbl:books3}, which show a modest increase in perplexity at the 2048k token length.

\textbf{Performance on established benchmarks within the original context window}. We assess LongRoPE-2048k models within the native context window on various tasks from the Hugging Face Open LLM Leaderboard~\cite{huggingfaceboard} in both zero-shot and few-shot configurations. Specifically, the benchmarks include ARC-Challenge~\cite{allenai:arc} with 25 shots, HellaSwag~\cite{zellers2019hellaswag} with 10 shots, MMLU~\cite{mmlu} with 5 shots, and TruthfulQA~\cite{lin2021truthfulqa} in a zero-shot setting.

Table~\ref{tbl:commonsense} illustrates that our models perform on par with the original benchmark, which was intended for a limited context window, and even surpass the original Mistral on TruthfulQA by an improvement of +0.5\%. Although LongRoPE-LLaMA2-2048k, which underwent fine-tuning at 256k, experiences marginally greater declines in performance, its results are still generally consistent with acceptable limits for the majority of evaluations.

\subsection{Ablation Results}

\textbf{Evaluating the effectiveness of the second positional interpolation}. Within our progressive extension framework, we employ our search-based method to perform a subsequent non-uniform positional interpolation on the extended, fine-tuned LLMs. To assess the impact of this approach, we experiment with our fine-tuned LLaMA2-256k model, expanding its context window to 512k, 1024k, and 2048k via PI and YaRN methods. The results, summarized in Table~\ref{tbl:secondsearch}, indicate that our non-uniform positional interpolation reliably maintains stable perplexity scores. Conversely, when applying PI and YaRN, we observe a rapid rise in perplexity as the extension factor grows.

\textbf{Recovery performance at reduced context lengths.} To address diminished effectiveness at shorter contexts, we adjust the RoPE scaling parameters for LongRoPE-2048k using our optimization algorithm. In this procedure, the upper limit for scale factors is reduced to minimize interpolation at context sizes of 4k and 8k tokens. Table~\ref{tbl:shortperformance} presents the perplexity outcomes for LongRoPE-LLaMA2-2048k evaluated on Proof-pile at both 4k and 8k tokens, together with the mean accuracy across LLM benchmarks. These findings indicate pronounced enhancements in performance at brief context lengths.

\textbf{Examination of the dual types of non-uniformities}. To assess the individual roles of the two forms of non-uniformity, we conduct ablation studies. Specifically, we perform two sets of experiments: (i) scaling LLaMA2-7B to shorter context lengths, namely 16k and 32k, via several approaches—PI, optimizing the RoPE dimension alone, and jointly tuning both non-uniformities; (ii) further increasing the length of our previously fine-tuned 256k LLaMA2 model up to 2048k, employing the same methodology. Perplexity measurements are carried out without any fine-tuning. Table~\ref{tbl:searchdimension} illustrates that modifying the non-uniformity in the RoPE dimension yields notably lower perplexity compared to the linear interpolation method used in PI. Adjusting token position non-uniformity leads to evident gains at 16k and 32k lengths, though similar benefits are absent at 2048k, likely attributable to the sheer length of the input. In these long contexts, retaining only the original tokens without any interpolation loses its efficacy. We propose to investigate this observation further in subsequent studies.

\section{Related Works}

Besides position interpolation methods, this section reviews other relevant approaches in the literature.

\textbf{Retrieval-based} methods incorporate an external memory component to store extensive previous context, along with retrieval mechanisms that fetch pertinent information during inference~\cite{longllama,wang2023augmenting,borgeaud2022improving}. These strategies often require substantial alterations to the underlying LLM architectures. In comparison, our method is more streamlined, introducing only slight changes to positional embeddings. Furthermore, our approach supports a broader range of long-context tasks, extending beyond retrieval to include long document summarization and few-shot learning.

\textbf{Attention-driven expansion of context windows}. In addition to modifying positional embeddings for interpolation, certain studies have enabled the extension of input context by leveraging the base LLM context window size through adjustments to attention mechanisms~\cite{han2023lminfinite, streamingllm,parallelwindow}. The central concept centers on alleviating the growth of attention computation arising from increased sequence positions, primarily through innovative attention masking approaches. Such methods work in tandem with positional interpolation strategies, offering complementary solutions.

\textbf{Fine-tuning based approaches} investigate strategies to adapt pre-trained LLMs with altered position embeddings, enabling them to handle extended contexts more effectively. Methods such as Code LLaMA~\cite{codellama}, LLaMA2 Long~\cite{xiong2023effective}, and ScaledRoPE~\cite{liu2023scaling} typically utilize a substantially larger RoPE base value and perform fine-tuning on the desired sequence length. By contrast, our approach provides greater adaptability to different context lengths and is capable of handling sequences surpassing 2M tokens. Recently, the high GPU resource demands inherent in fine-tuning for extremely long contexts (e.g., over 128k tokens) have led to the introduction of methods like LongLoRA~\cite{longlora} and PoSE~\cite{zhu2023pose}, which aim to reduce the associated computational burden. Our method operates independently and can be combined with these efficiency-focused fine-tuning techniques.

\section{Conclusion}

This paper introduces LongRoPE, a technique that significantly increases the context window of LLMs to an extraordinary 2048k, all while preserving their performance within standard, shorter context ranges. Our approach leverages two distinct non-uniformities in RoPE positional embeddings and applies an efficient evolutionary search. This strategy yields dual advantages: it supplies robust initialization for subsequent fine-tuning and allows for an 8$\times$ expansion of the context window even in the absence of fine-tuning. Building upon these insights, we propose a progressive extension process that employs LLMs fine-tuned at 256k-length contexts, ultimately achieving a 2048k window without the need for additional fine-tuning. Comprehensive experimental evaluation demonstrates the efficacy of LongRoPE. We anticipate that LongRoPE-2048k will unlock novel applications reliant on extensive context and foster new directions in this area of research.

\section*{Broader Impacts} The research described in this paper aims to contribute to progress in Machine Learning. While our work may carry a range of societal implications, we do not identify any particular outcomes that warrant explicit mention in this section.

	\balance

\end{document}